\documentclass[11pt,letterpaper]{article}
% Reusable style. Build each .tex from its own directory using scripts/build.sh.
\usepackage[T1]{fontenc}
\usepackage[utf8]{inputenc}
\usepackage{lmodern}
\usepackage[margin=0.86in,headheight=15pt]{geometry}
\usepackage{amsmath,amssymb,amsthm,mathtools}
\usepackage{booktabs,longtable,tabularx,array,multicol}
\usepackage[table]{xcolor}
\usepackage{microtype}
\usepackage{enumitem}
\usepackage{fancyhdr}
\usepackage{titlesec}
\usepackage{hyperref}
\usepackage{listings}
\usepackage[most]{tcolorbox}
\usepackage{needspace}
\definecolor{aeNavy}{HTML}{172D43}
\definecolor{aeBlue}{HTML}{176B90}
\definecolor{aeTeal}{HTML}{148477}
\definecolor{aePale}{HTML}{F2F7F9}
\definecolor{aeGray}{HTML}{5C6771}
\hypersetup{colorlinks=true,linkcolor=aeBlue,urlcolor=aeBlue,pdfauthor={Algorithm Engineering Documentation Suite},pdfsubject={Independent educational materials based on the supplied TAOCP volume topics}}
\titleformat{\section}{\Large\bfseries\color{aeNavy}}{\thesection}{0.6em}{}
\titleformat{\subsection}{\large\bfseries\color{aeBlue}}{\thesubsection}{0.6em}{}
\titleformat{\subsubsection}{\normalsize\bfseries\color{aeNavy}}{\thesubsubsection}{0.6em}{}
\setlist[itemize]{leftmargin=*,topsep=3pt,itemsep=1pt}
\setlist[enumerate]{leftmargin=*,topsep=3pt,itemsep=3pt}
\setlength{\parindent}{0pt}\setlength{\parskip}{5pt}
\pagestyle{fancy}\fancyhf{}\fancyhead[L]{\footnotesize\color{aeGray}ALGORITHM ENGINEERING}\fancyhead[R]{\footnotesize\color{aeGray}DOCUMENTATION SUITE}\fancyfoot[L]{\footnotesize\color{aeGray}Independent study and engineering practice}\fancyfoot[R]{\footnotesize\thepage}
\renewcommand{\headrulewidth}{0.35pt}
\newtcolorbox{engineeringnote}[1][]{breakable,colback=aePale,colframe=aeBlue,boxrule=0.5pt,arc=1mm,fonttitle=\bfseries,title={Engineering note},#1}
\newtcolorbox{evidencebox}[1][]{breakable,colback=white,colframe=aeTeal,boxrule=0.6pt,arc=1mm,fonttitle=\bfseries,title={Evidence and deliverables},#1}
\lstdefinestyle{algo}{basicstyle=\ttfamily\footnotesize,breaklines=true,frame=single,rulecolor=\color{aeGray!40},backgroundcolor=\color{aePale},showstringspaces=false,tabsize=2,columns=fullflexible,keepspaces=true,keywordstyle=\color{aeBlue}\bfseries,commentstyle=\color{aeTeal}}
\lstset{style=algo}
\newcommand{\DocTitle}[3]{%
\hypersetup{pageanchor=false}\begin{titlepage}\vspace*{1.2cm}
{\color{aeTeal}\rule{\linewidth}{2.3pt}}\par\vspace{1.15cm}
{\fontsize{24}{30}\selectfont\bfseries\color{aeNavy}\raggedright #1\par}\vspace{0.45cm}
{\Large\color{aeBlue}#2\par}\vspace{1.3cm}
{\large\bfseries #3\par}\vspace{0.35cm}
{\color{aeGray}Edition 1.0 \quad | \quad 29 October 2026\par}\vfill
{\color{aeGray}\small Original educational material aligned to the chapter and section topics in the five supplied volumes of Donald E. Knuth's \emph{The Art of Computer Programming}. This document is not an authorized reproduction or substitute for the original books.}\par\vspace{0.4cm}
{\color{aeTeal}\rule{\linewidth}{0.75pt}}
\end{titlepage}\hypersetup{pageanchor=true}\pagenumbering{roman}}
\newcommand{\SourceNote}[1]{\par\smallskip{\footnotesize\color{aeGray}\textbf{Source locator:} #1}\par\smallskip}
\newcommand{\Outcome}[1]{\par\textbf{Competency target.} #1\par}
\newcommand{\Checkpoint}[1]{\par\textbf{Checkpoint.} #1\par}
% Intentionally simple printable writing lines: call \AnswerLines{number}.
\newcount\alcount
\newcommand{\AnswerLines}[1]{\par\alcount=0\loop\ifnum\alcount<#1\par\noindent\rule{\linewidth}{0.3pt}\vspace{0.33cm}\advance\alcount by1\repeat\par}
\newcommand{\FormField}[2][0.62\linewidth]{\noindent\textbf{#2}\hspace{0.65em}\rule{#1}{0.4pt}\par\vspace{0.18cm}}
\newcommand{\ReviewTable}{\begin{tabularx}{\linewidth}{@{}p{0.23\linewidth}X p{0.2\linewidth}@{}}\toprule Criterion & Recorded evidence & Rating (0--4)\\\midrule Correctness & & \\[1em] Proof quality & & \\[1em] Complexity & & \\[1em] Testing & & \\[1em] Reproducibility & & \\[1em] Communication & & \\[1em]\bottomrule\end{tabularx}}

\begin{document}
\DocTitle{Randomness and Statistical Evaluation Workbook}{Worked models, original problems and writing space}{Mathematical analysis workbook}
\tableofcontents
\clearpage
\pagenumbering{arabic}
\section{Source alignment}\SourceNote{Vol. 2, Chapter 3.}

\section{Worked model: modular recurrence}
For modulus $m=9$, multiplier $a=4$, increment $c=1$, seed $x_0=0$, successive states begin $0,1,5,3,4,8,6,7,2,0$. A complete return to the initial state occurs after nine steps for this chosen example. This single case establishes neither quality nor cryptographic security.
\section{Worked model: testing a uniform categorical hypothesis}
For $k$ equiprobable categories, expected count in $N$ samples is $E_i=N/k$ and $\chi^2=\sum_i(O_i-E_i)^2/E_i$. Its interpretation depends on assumptions, sample size and an appropriate reference distribution. A non-rejection does not prove that observations are independent.

\section{Workbook protocol}
Write assumptions before derivations. Distinguish claims to prove, counterexamples to find and empirical questions to test. Use separate paper if a proof exceeds the available lines. Each response should end with a testable implication or explicit limitation.

\subsection{Exercise 1}\noindent Find the period of a selected LCG for modulus 16, checking all seeds.\AnswerLines{3}
\subsection{Exercise 2}\noindent Design a histogram test and a consecutive-pair test that detect different defects.\AnswerLines{3}
\subsection{Exercise 3}\noindent Write a reproducible Fisher--Yates experiment over four symbols.\AnswerLines{3}
\subsection{Exercise 4}\noindent Calculate the expected frequency of each permutation in 24,000 unbiased shuffles.\AnswerLines{3}
\subsection{Exercise 5}\noindent Explain when chi-squared expected counts become too small.\AnswerLines{3}
\subsection{Exercise 6}\noindent Distinguish seed repeatability from unpredictable entropy.\AnswerLines{3}
\subsection{Exercise 7}\noindent Construct an example with uniform marginal frequencies but strong serial dependence.\AnswerLines{3}
\subsection{Exercise 8}\noindent Describe how multiple testing changes false-positive interpretation.\AnswerLines{3}
\subsection{Exercise 9}\noindent Propose a rejection sampling routine for six equiprobable outcomes from random bits.\AnswerLines{3}
\subsection{Exercise 10}\noindent Identify a threat model for which a simulation PRNG would be disallowed.\AnswerLines{3}
\section{Self-review}
\begin{itemize}\item Every symbol has a domain.\item All bounding inequalities declare their valid range.\item Inputs and edge cases are explicit.\item Statistical and asymptotic claims state their model.\item Code results are checked against an independent calculation.\end{itemize}
\end{document}
